%% ma: L10-B12-0056
%% lop: 10
%% chuong: 5
%% bai: 12
%% dang: 10.12.5
%% loai: TN4
%% muc: TH
%% dap_an: "B"
%% nguon: "chuyen-de-cac-so-dac-trung-cua-mau-so-lieu-khong-ghep-nhom-toan-10-chuong-trinh-moi.pdf (GV. Phan Nhật Linh) (P01), Bài 01 Phần I Câu 40"
%% kien_thuc: "Lớp 10 Bài 12, đổi đơn vị thời gian, kí hiệu khoa học"
%% trang_thai: chinh-thuc
%% ghi_chu: "Đáp án gốc 9,4608.10^9 sai (quên nhân 7 và nhân 1000); đáp án đúng 6,62256.10^13 km, đã thay bộ phương án. Tự duyệt theo yêu cầu giáo viên 22/10/2026."
\begin{ex}
Một máy bay giả định có tốc độ gấp bảy lần tốc độ ánh sáng. Biết tốc độ ánh sáng là $300$ nghìn km/s và một năm có $365$ ngày. Quãng đường máy bay đó bay được trong một năm, viết dưới dạng kí hiệu khoa học, là
\choice{$9{,}4608\cdot10^{12}$ km}{\True $6{,}62256\cdot10^{13}$ km}{$6{,}62256\cdot10^{10}$ km}{$9{,}4608\cdot10^{9}$ km}
\loigiai{Một năm có $365\cdot24\cdot60\cdot60=31\,536\,000$ giây.

Quãng đường: $7\cdot300\,000\cdot31\,536\,000=6{,}62256\cdot10^{13}$ (km). Chọn B.}
\end{ex}
